\documentclass{article}
\usepackage[T1]{fontenc}
\usepackage{lmodern}
\usepackage{graphicx}
\usepackage{amsmath,amssymb}
\usepackage[a4paper,margin=2cm]{geometry}
\usepackage[absolute,overlay]{textpos}
\setlength{\TPHorizModule}{1mm}
\setlength{\TPVertModule}{1mm}
\usepackage[indonesian]{babel}
\usepackage{hyperref}
\setlength{\emergencystretch}{2em}
% unit: Halaman 23

\begin{document}

% === Halaman 23 (llm) ===

\section*{Kriptanalisis Knapsack}

\begin{itemize}
    \item Sayangnya, algoritma \textit{knapsack} dinyatakan sudah tidak aman, karena \textit{knapsack} dapat dipecahkan oleh pasangan kriptografer Adi Shamir.
    \item Shamir merumuskan transformasi yang memungkinkan mereka merekonstruksi \textit{superincreasing knapsack} dari \textit{normal knapsack} dalam waktu polynomial.
    \item Baca makalahnya: \textit{A polynomial time algorithm for breaking the basic Merkle-Hellman cryptosystem}

    (\textbf{Published in:} \href{https://doi.org/10.1145/800014.800015}{23rd Annual Symposium on Foundations of Computer Science (sfcs 1982)})
\end{itemize}

\end{document}
